\chapter{Rheumatology}

\medterm{Achilles Tendinopathy}-- Degeneration/inflammation of Achilles tendon. Symptoms: posterior heel pain, thickening, stiffness. Risk factors: running, improper footwear, tight calf muscles. Treatment: eccentric exercises, heel raises, orthotics, shockwave therapy, surgery if refractory.

\medterm{Adult-Onset Still Disease}-- AOSD: a rare systemic autoinflammatory disorder characterized by daily high-spiking fevers, evanescent salmon-pink rash, arthritis/arthralgia, sore throat, lymphadenopathy, hepatosplenomegaly, and markedly elevated ferritin (with low glycosylated ferritin), ESR/CRP. Complications: macrophage activation syndrome, pericarditis, DIC. Diagnosis: Yamaguchi criteria; exclude infection, malignancy, other rheumatic disease. Management: NSAIDs, corticosteroids (first-line for moderate-severe), methotrexate, IL-1 inhibitors (anakinra, canakinumab), IL-6 inhibitor (tocilizumab), and TNF inhibitors.

\medterm{ANCA-Associated Vasculitis}-- A group of small-vessel vasculitides associated with antineutrophil cytoplasmic antibodies. Includes granulomatosis with polyangiitis (c-ANCA/PR3), microscopic polyangiitis (p-ANCA/MPO), and eosinophilic granulomatosis with polyangiitis (p-ANCA/MPO with asthma/eosinophilia). Treatment: glucocorticoids plus rituximab or cyclophosphamide for induction, rituximab or azathioprine for maintenance.

\medterm{Ankylosing Spondylitis}-- A chronic inflammatory disease primarily affecting the axial skeleton, causing fusion of vertebrae (bamboo spine) and sacroiliitis. Associated with HLA-B27 in 90% of patients. Symptoms: inflammatory back pain (worse with rest, better with exercise), morning stiffness, limited spinal mobility. Complications: uveitis, aortitis, apical pulmonary fibrosis, amyloidosis. Treatment: NSAIDs first-line, TNF inhibitors or IL-17 inhibitors for refractory disease, exercise and physical therapy essential.

\medterm{Anterior Cruciate Ligament Tear}-- ACL tear: common knee injury, often non-contact pivoting sports. Symptoms: popping sound, swelling, instability, inability to continue play. Diagnosis: Lachman test, anterior drawer, MRI. Treatment: surgical reconstruction (autograft/allograft), rehabilitation.

\medterm{Anterior/Posterior Drawer Test}-- Knee ligament testing. Anterior (ACL), posterior (PCL).
\medterm{Anti-CCP}-- Anti-cyclic citrullinated peptide antibody: an autoantibody against citrullinated proteins, highly specific (90-95 percent) for rheumatoid arthritis and predictive of erosive disease. Positive anti-CCP is a key component of ACR/EULAR RA classification; may be present years before clinical onset. Used with rheumatoid factor; anti-CCP-positive, RF-negative and vice versa occur. Useful in early/inflammatory arthritis to identify patients needing aggressive DMARD therapy. Does not correlate directly with disease activity over time.

\medterm{Antinuclear Antibody}-- ANA: autoantibodies against nuclear components (DNA, histones, ribonucleoproteins, etc.), detected by immunofluorescence (HEp-2 cells), reported with pattern (homogeneous, speckled, nucleolar, centromere) and titer. ANA is a screening test for systemic autoimmune disease: high positive in SLE (99 percent), also in drug-induced lupus, Sjogren, scleroderma, MCTD, RA, and in healthy individuals (low titers). Positive ANA requires specific antibody testing based on pattern/suspicion (anti-dsDNA, anti-Smith—SLE; anti-SSA/SSB—Sjogren; anti-Scl-70, centromere—scleroderma; anti-RNP—MCTD). ANA is not diagnostic alone.

\medterm{Antiphospholipid Antibodies}-- Include lupus anticoagulant (prolongs phospholipid-dependent coagulation tests), anticardiolipin antibodies (IgG/IgM), and anti-beta-2-glycoprotein I antibodies. Detected on >=2 occasions 12 weeks apart. Clinical criteria: vascular thrombosis or pregnancy morbidity. Can be transient (infection, drugs) or persistent.

\medterm{Antiphospholipid Syndrome}-- A thrombophilic disorder with arterial/venous thrombosis and/or pregnancy morbidity, with persistent antiphospholipid antibodies (lupus anticoagulant, anticardiolipin, anti-beta-2-glycoprotein I). Can be primary or secondary to SLE. Management: anticoagulation with warfarin (target INR 2-3 for venous thrombosis, 2-3 or 3-4 for arterial); antiplatelet therapy for pregnancy.

\medterm{Antisynthetase Syndrome}-- Defined by the presence of anti-aminoacyl-tRNA synthetase antibodies (anti-Jo-1 most common, 75%; anti-PL-7, anti-PL-12, anti-EJ, anti-OJ, anti-KS, anti-Zo, anti-Ha, anti-YRS, anti-FRS, anti-TRS, anti-ARS) with clinical triad: interstitial lung disease (most common, often progressive), inflammatory myopathy (DM/PM), and inflammatory arthritis (non-erosive). Also: mechanic's hands (hyperkeratotic fissuring of radial fingers), Raynaud, fever. Treatment: glucocorticoids + rituximab or mycophenolate for ILD and myositis; calcineurin inhibitors (tacrolimus, cyclosporine) for refractory disease.

\medterm{Aphthous Ulcer}-- Recurrent, painful, round or oval ulcers of the oral mucosa, commonly called canker sores. See dentistry.

\medterm{Arthritis}-- Inflammation of one or more joints, characterized by joint pain, swelling, warmth, and stiffness. Arthritis is classified into several major categories: osteoarthritis (degenerative, the most common form, characterized by articular cartilage loss, subchondral sclerosis, and osteophyte formation, typically affecting weight-bearing joints such as knees, hips, and spine as well as hands); rheumatoid arthritis (autoimmune, symmetric polyarthritis affecting small joints of hands and feet with pannus formation causing marginal erosions and periarticular osteopenia); spondyloarthritis (including ankylosing spondylitis, psoriatic arthritis, reactive arthritis, and enteropathic arthritis, typically presenting with inflammatory back pain, asymmetric oligoarthritis, enthesitis, and association with HLA-B27); crystal arthropathies (gout and calcium pyrophosphate deposition disease); septic arthritis (bacterial infection of the joint, a medical emergency requiring urgent drainage and antibiotics); and other arthritides associated with connective tissue diseases, infection, or haemochromatosis.

\medterm{Autoimmune Hepatitis}-- Chronic hepatitis caused by immune-mediated destruction of hepatocytes characterized by interface hepatitis with lymphoplasmacytic infiltrate, elevated transaminases, hypergammaglobulinemia, and autoantibodies. See \hyperlink{Gastroenterology}{gastroenterology}.

\medterm{Avascular Necrosis}-- See vascular surgery. Osteonecrosis of bone due to interrupted blood supply. Common sites: femoral head, humeral head, talus. Risk factors: steroids, alcohol, trauma, sickle cell, SLE. Treatment: core decompression, joint replacement.

\medterm{Baker Cyst}-- Popliteal cyst behind knee, often associated with knee pathology (OA, meniscal tear, RA). May rupture causing calf pain/swelling (pseudothrombophlebitis). Diagnosis: ultrasound/MRI. Treatment: treat underlying condition, aspiration if symptomatic.

\medterm{Bankart Lesion}-- Injury to inferior glenohumeral ligament/labrum from shoulder dislocation. Associated with recurrent instability. Treatment: surgical repair (arthroscopic/open).

\medterm{Bechterew Disease}-- Historical eponym for ankylosing spondylitis (also known as Marie-Strümpell disease), describing the chronic inflammatory arthritis of the axial skeleton first reported by Vladimir Bekhterev (Russian, 1892) and Adolf Strümpell (German, 1884). Characterized by inflammatory back pain, sacroiliitis, syndesmophyte formation, eventual "bamboo spine" with ankylosis, and association with HLA-B27 in approximately $90\%$ of cases. Modern usage prefers ankylosing spondylitis / axial spondyloarthritis; "Bechterew disease" remains in some European (Russian, German) medical literature. See \hyperlink{Ankylosing Spondylitis}{Ankylosing Spondylitis}.

\medterm{Biologic Therapy}-- Treatment with biologic agents targeting specific immune pathways, used in autoimmune and inflammatory disease. Types: TNF inhibitors (infliximab, etanercept, adalimumab), IL-6 inhibitors (tocilizumab), IL-1 inhibitors (anakinra), IL-17/IL-23 inhibitors (secukinumab, ustekinumab), B-cell depletion (rituximab), T-cell costimulation blockade (abatacept), JAK inhibitors (tofacitinib—small molecule), and others. Indications: rheumatoid arthritis, psoriatic arthritis, ankylosing spondylitis, SLE, vasculitis, IBD, psoriasis. Benefits: targeted efficacy with steroid-sparing. Risks: infection (TB, opportunistic), malignancy (rare), immunogenicity (anti-drug antibodies), and injection reactions. Requires screening and monitoring.

\medterm{Bursitis}-- Inflammation of the bursa (fluid-filled sac reducing friction between tissues). Common sites: subacromial (shoulder), olecranon (elbow), prepatellar (knee), trochanteric (hip). Causes: repetitive stress, trauma, infection, crystal deposition (gout, pseudogout), autoimmune disease. Treatment: rest, ice, NSAIDs, corticosteroid injection, aspiration if septic.


\medterm{Calcium Pyrophosphate Deposition Disease}-- A crystal arthropathy caused by deposition of calcium pyrophosphate dihydrate crystals in articular cartilage (chondrocalcinosis). Presents with acute pseudogout (most common): acute monoarticular inflammation of the knee or wrist, mimicking gout but less painful. Chronic CPPD: degenerative arthritis-like with OA pattern (knee, wrist, MCP joints). Pseudo-rheumatoid: symmetric polyarticular inflammation. Associated with aging, OA, hemochromatosis, hyperparathyroidism, hypomagnesemia, hypothyroidism. Diagnosis: joint aspiration shows weakly positively birefringent rhomboid crystals. Treatment: NSAIDs, colchicine, intra-articular corticosteroids for acute flares; low-dose colchicine for prophylaxis.

\medterm{Carpal Tunnel Syndrome}-- Median nerve compression at wrist. Symptoms: numbness/tingling in thumb, index, middle, radial half of ring finger. Night symptoms common. Diagnosis: clinical, nerve conduction studies. Treatment: splinting, corticosteroid injection, surgical decompression.

\medterm{Chondrocalcinosis}-- Calcification of cartilage visible on X-ray. Associated with CPPD disease. Found in knees, wrists, pubic symphysis. May be incidental or symptomatic.

\medterm{Chronic Fatigue Syndrome}-- ME/CFS: debilitating condition characterized by severe fatigue lasting >6 months, post-exertional malaise, unrefreshing sleep, cognitive impairment, orthostatic intolerance. Diagnosis: clinical, exclude other causes. Treatment: symptom management, pacing, graded exercise therapy controversial.

\medterm{Collateral Ligament Injury}-- MCL (medial) or LCL (lateral) knee ligament injury. MCL more common, often from valgus stress. Treatment: braces, physical therapy, surgical repair if severe.

\medterm{Complex Regional Pain Syndrome}-- CRPS: chronic pain condition, usually after trauma/surgery. Type I (no nerve injury), Type II (nerve injury). Symptoms: pain disproportionate to injury, allodynia, swelling, color/temperature changes, decreased range of motion. Treatment: physical therapy, medications (bisphosphonates, gabapentin, corticosteroids), sympathetic blocks.

\medterm{Connective Tissue Disease}-- Autoimmune diseases affecting connective tissue. Examples: SLE, scleroderma, RA, Sjogren syndrome, mixed connective tissue disease. Overlap syndromes common.

\medterm{CREST Syndrome}-- Limited cutaneous form of systemic sclerosis (scleroderma), characterized by the eponymic features: Calcinosis cutis (multiple subcutaneous calcium deposits, often in extremities — can cause pain, inflammation, ulceration), Raynaud phenomenon (usually long-standing; often precedes other features by years), Esophageal dysmotility (lower two-thirds smooth muscle, causes dysphagia, reflux, strictures), Sclerodactyly (thickening and tightening of skin distal to MCP joints — fingers become tapered, immobile), Telangiectasias (mat-like, on face, lips, palms, finger pads). Anti-centromere antibodies (anti-CENP-B) positive in 70-80\% of cases (highly specific for limited scleroderma; less specific in diffuse disease — anti-Scl-70 (anti-topoisomerase I) is more associated with diffuse). Clinical course: relatively indolent compared with diffuse systemic sclerosis (skin thickening limited to distal extremities and face, lung involvement later and slower; pulmonary arterial hypertension (PAH) in 10-15\%), unlike diffuse disease where early ILD and renal crisis predominate. Diagnosis: ACR/EULAR 2013 classification (skin thickening of fingers extending proximal to MCP joints suffices); nailfold capillaroscopy; antibody testing. Management: Raynaud — calcium channel blockers (amlodipine, nifedipine), phosphodiesterase type 5 inhibitors (sildenafil), prostacyclin analogs, digital sympathectomy for severe digital ischemia; PPI for reflux; vasodilator therapy for PAH (bosentan, ambrisentan, macitentan, sildenafil, riociguat); immunosuppression generally limited in limited cutaneous SSc; periodic surveillance echocardiography, PFTs, ECG.

\medterm{Cryopyrin-Associated Periodic Syndromes}-- Autosomal dominant autoinflammatory disorders (NLRP3 mutations, gain-of-function) causing constitutive inflammasome activation and IL-1beta overproduction. Spectrum: familial cold autoinflammatory syndrome (FCAS, cold-induced urticaria, fever, conjunctivitis), Muckle-Wells syndrome (MWS, progressive sensorineural hearing loss, amyloidosis risk), neonatal-onset multisystem inflammatory disease (NOMID/CINCA, severe CNS involvement, bony overgrowth, chronic meningitis, cognitive impairment). Diagnosis is clinical with genetic confirmation. Treatment: IL-1 inhibitors (anakinra, canakinumab, rimonacept) are dramatically effective, suppressing inflammation and preventing disease progression.

\medterm{Crystal Arthropathy}-- Joint disease caused by crystal deposition. Gout (monosodium urate) and pseudogout (calcium pyrophosphate). Diagnosis by joint aspiration and polarized microscopy.

\medterm{Cubital Tunnel Syndrome}-- Ulnar nerve compression at elbow. Symptoms: numbness/tingling in little finger, medial hand. Elbow flexion exacerbates. Treatment: elbow padding, splinting, surgical decompression/transposition.

\medterm{Dactylitis}-- Inflammation of an entire digit (finger or toe), producing a diffuse, fusiform swelling often described as sausage-shaped. Most commonly associated with psoriatic arthritis and reactive arthritis (enthesitis-related), but also occurs in sickle cell disease (hand-foot syndrome in young children, caused by dactylitis from vaso-occlusive crises in the small bones of the hands and feet) and tuberculous dactylitis (spina ventosa--bony expansion from granulomatous infection). Diagnosis based on clinical appearance and underlying condition.

\medterm{De Quervain Tenosynovitis}-- Tenosynovitis of first dorsal compartment tendons (abductor pollicis longus, extensor pollicis brevis). Finkelstein test positive. Treatment: splinting, corticosteroid injection, surgical release.

\medterm{Dermatomyositis}-- An inflammatory myopathy with characteristic skin manifestations: Gottron papules (over knuckles), heliotrope rash (purple periorbital), shawl/V sign erythema, and proximal muscle weakness; associated with interstitial lung disease, and in adults an increased risk of malignancy (ovary, lung, GI). Antibodies: anti-Mi-2, anti-Jo-1 (antisynthetase with ILD, mechanic's hands), anti-MDA5 (rapidly progressive ILD). Diagnosis: elevated CK, EMG, muscle biopsy (perifascicular atrophy), skin biopsy, MRI. Management: corticosteroids, methotrexate/azathioprine, IVIG, rituximab; cancer screening at diagnosis. Juvenile dermatomyositis variant in children.

\medterm{Drug-Induced Autoimmunity}-- Drug-induced lupus (DIL): procainamide, hydralazine, isoniazid, minocycline, TNF inhibitors, anti-TNF; anti-histone antibodies, positive ANA, negative anti-dsDNA/anti-Sm; symptoms resolve after drug cessation. Drug-induced ANCA vasculitis: propylthiouracil, hydralazine, minocycline, allopurinol; MPO-ANCA positive, pauci-immune GN, pulmonary hemorrhage. Drug-induced subacute cutaneous lupus: hydrochlorothiazide, calcium channel blockers, TNF inhibitors; anti-Ro/SSA positive, annular/papulosquamous photosensitive rash. Diagnosis is clinical with drug history and serology. Management is withdrawal of the offending drug with supportive care; symptoms typically resolve over weeks to months.

\medterm{Dupuytren Contracture}-- Progressive fibrosis of palmar fascia causing flexion contracture of fingers (ring, little). More common in men, Northern European descent, diabetes, epilepsy. Treatment: observation, collagenase injection, fasciectomy, needle aponeurotomy.

\medterm{Ehlers-Danlos Syndrome}-- Group of connective tissue disorders affecting collagen. Types: hypermobile, classical, vascular, arthrochalasic, kyphoscoliotic, dermatosparaxis. Features: joint hypermobility, skin hyperextensibility, tissue fragility. Vascular type: arterial/intestinal/uterine rupture risk. Diagnosis: clinical criteria, genetic testing. Management: multidisciplinary, avoid contact sports, genetic counseling.

\medterm{Enteropathic Arthritis}-- Spondyloarthritis associated with inflammatory bowel disease (Crohn disease, ulcerative colitis), affecting 10-20% of IBD patients. Two patterns: type 1 (pauciarticular, acute, self-limited, parallels IBD activity, large joints), type 2 (polyarticular, chronic, independent of IBD activity, symmetric small joints, resembles RA but RF-negative). Axial involvement (sacroiliitis) in 10-20%, often asymptomatic. Treatment: optimize IBD control (TNF inhibitors, vedolizumab for gut-selective therapy), NSAIDs for symptomatic relief (though they may exacerbate IBD), sulfasalazine, methotrexate, and rarely surgery for refractory arthritis.

\medterm{Enthesitis}-- Inflammation at the site where a tendon, ligament, or joint capsule inserts into bone (enthesis), a hallmark of spondyloarthritis (ankylosing spondylitis, psoriatic arthritis, reactive arthritis, IBD-associated). Affects Achilles tendon, plantar fascia, patellar, and costochondral junctions. Features: localized pain and tenderness at insertion sites, morning stiffness, sometimes swelling; plantar fasciitis and Achilles enthesitis are common. Diagnosis: clinical, ultrasound/MRI (entheseal thickening, vascularity, erosions). Management: NSAIDs, physiotherapy, and in spondyloarthritis—biologics (TNF inhibitors, IL-17 inhibitors) for refractory enthesitis.

\medterm{Eosinophilic Granulomatosis with Polyangiitis}-- EGPA (formerly Churg-Strauss syndrome): an ANCA-associated vasculitis characterized by asthma, peripheral and tissue eosinophilia, and necrotizing vasculitis. Features: late-onset asthma, allergic rhinitis, eosinophilic pneumonia, mononeuritis multiplex, skin lesions, cardiac involvement (myocarditis—major cause of death), and glomerulonephritis (less common); p-ANCA (anti-MPO) in ~40-50 percent. Diagnosis: clinical, eosinophilia (>10 percent/1500), biopsy (eosinophilic vasculitis), ANCA. Management: corticosteroids (mainstay), cyclophosphamide or rituximab for severe/organ-threatening; treat asthma and eosinophilia; relapse common.

\medterm{Erythrocyte Sedimentation Rate}-- ESR: a nonspecific inflammatory marker measuring the rate of red blood cell sedimentation in a vertical tube over 1 hour (mm/h), elevated by inflammation, infection, malignancy, and anemia. Extremely elevated (>100 mm/h): myeloma, giant cell arteritis/polymyalgia rheumatica, infection, SLE, malignancy. Used with CRP to monitor inflammation (RA, GCA) and in diagnosis (temporal arteritis). Affected by age, sex, anemia (increases), polycythemia (decreases), fibrinogen, and pregnancy. Not specific—combination with clinical findings and CRP is standard.

\medterm{Familial Mediterranean Fever}-- An autosomal recessive autoinflammatory disorder (MEFV gene mutations, pyrin/inflammasome dysregulation) causing recurrent self-limited febrile attacks with serositis (peritonitis, pleuritis, synovitis), erysipelas-like erythema, lasting 12-72 hours. Predominantly affects Mediterranean populations (Armenians, Turks, Arabs, Sephardic Jews). Complication: AA amyloidosis (renal, hepatic, splenic) from chronic inflammation if untreated. Diagnosis: Tel-Hashomer criteria: major (typical attacks with peritonitis, pleuritis, synovitis, or erysipelas-like erythema) plus minor (incomplete attacks). Treatment: colchicine is highly effective for preventing attacks and AA amyloidosis; IL-1 inhibitors (anakinra, canakinumab) for colchicine-resistant cases.

\medterm{Fasciitis}-- Inflammation of fascia. Examples: plantar fasciitis (heel pain), eosinophilic fasciitis (Shulman disease—skin thickening, eosinophilia). Treatment: NSAIDs, corticosteroids, physical therapy.

\medterm{Felty Syndrome}-- Triad of rheumatoid arthritis, splenomegaly, and neutropenia (ANC <2000). Complication of long-standing, severe, seropositive RA. Increased infection risk. Treatment: DMARDs (methotrexate, TNF inhibitors), G-CSF for neutropenia, splenectomy in refractory cases.

\medterm{Fibromyalgia}-- A chronic centralized pain syndrome characterized by widespread musculoskeletal pain (>3 months), fatigue, sleep disturbance, cognitive dysfunction (fibro fog), and somatic symptoms, affecting 2-4% of population (female:male 9:1). Pathophysiology: central sensitization, altered pain processing, neurotransmitter dysregulation (serotonin, norepinephrine, dopamine, substance P). 2016 ACR criteria: widespread pain index (WPI) >=7 and symptom severity scale (SSS) >=5 OR WPI 4-6 and SSS >=9, plus generalized pain in at least 4 of 5 regions, symptoms present at similar level for at least 3 months, and absence of another disorder explaining the pain. Management: non-pharmacological (exercise, cognitive behavioural therapy, sleep hygiene, patient education) combined with pharmacotherapy (amitriptyline, duloxetine, milnacipran, pregabalin, gabapentin). Multimodal, individualised treatment is most effective.

\medterm{Finkelstein Test}-- Wrist deviation test for De Quervain tenosynovitis. Pain over radial styloid.

\medterm{Frozen Shoulder}-- Adhesive capsulitis: progressive shoulder pain and stiffness. Stages: freezing (pain), frozen (stiffness), thawing (recovery). Associated with diabetes, thyroid disease, immobilization. Treatment: physical therapy, corticosteroid injection, manipulation under anesthesia, arthroscopic capsular release.

\medterm{Giant Cell Arteritis}-- A large-vessel vasculitis of adults >50 (peak 70-80) affecting cranial branches of carotid arteries, with risk of permanent vision loss from anterior ischemic optic neuropathy. Clinical: new-onset headache (temporal), scalp tenderness, jaw claudication, visual disturbances (amaurosis fugax, diplopia), polymyalgia rheumatica symptoms. Diagnosis: temporal artery biopsy (gold standard, skip lesions reduce sensitivity), ultrasound (halo sign), PET-CT (large-vessel involvement (aorta, subclavian). Treatment: high-dose glucocorticoids (prednisolone 40-60 mg/day) promptly to prevent blindness; tocilizumab (IL-6 inhibitor) as steroid-sparing agent; low-dose aspirin for thromboprophylaxis. Methotrexate, azathioprine as alternatives.

\medterm{Gout}-- A crystal deposition disease caused by monosodium urate (MSU) crystal formation in joints and tissues due to chronic hyperuricemia (serum urate >6.8 mg/dL). Pathophysiology: urate overproduction (10%) or underexcretion (90%); crystals activate NLRP3 inflammasome -> IL-1beta release -> neutrophilic inflammation. Clinical: acute monoarticular arthritis (1st MTP podagra 50% first attack), intense pain, erythema, swelling; tophi in chronic disease; urate nephrolithiasis. Diagnosis: synovial fluid aspiration showing negatively birefringent needle-shaped crystals under polarised light microscopy, elevated serum urate. Treatment: acute attack management with NSAIDs, colchicine, or corticosteroids; long-term urate-lowering therapy (allopurinol, febuxostat, probenecid, uricosurics) targeting serum urate <6 mg/dL (5 mg/dL if tophi present). Prophylaxis with NSAIDs or colchicine during initial urate-lowering. Lifestyle modification: avoid purine-rich foods (red meat, shellfish, organ meats), reduce alcohol (especially beer), maintain hydration.

\medterm{Granulomatosis with Polyangiitis}-- GPA (formerly Wegener granulomatosis): an ANCA-associated small-vessel necrotizing vasculitis affecting the upper and lower respiratory tract and kidneys. Features: sinonasal disease (crusting, epistaxis, saddle-nose deformity), pulmonary nodules/cavitation, alveolar hemorrhage, glomerulonephritis (RPGN), and constitutional symptoms; c-ANCA (anti-PR3) positive in most. Diagnosis: clinical, ANCA serology, biopsy (necrotizing granulomas, vasculitis). Management: induction with rituximab or cyclophosphamide + high-dose corticosteroids, then maintenance (rituximab, azathioprine); treat renal failure aggressively; relapses require re-treatment. Untreated: rapidly fatal.

\medterm{Henoch-Schonlein Purpura}-- IgA vasculitis: palpable purpura on buttocks/legs, abdominal pain, arthralgia, renal involvement (IgA nephropathy), most common in children. Self-limited in most cases; corticosteroids for severe GI/renal involvement.

\medterm{Hill-Sachs Lesion}-- Impaction fracture of posterolateral humeral head from anterior shoulder dislocation. Associated with Bankart lesion.

\medterm{HLA-B27}-- Human leukocyte antigen B27, a class I MHC molecule strongly associated with spondyloarthropathies. Present in 90% of ankylosing spondylitis patients, 60-80% of reactive arthritis, and associated with acute anterior uveitis, psoriatic arthritis, and IBD-associated arthritis. Not diagnostic alone—present in 8% of healthy populations. Used in conjunction with clinical features.

\medterm{Hydroxychloroquine}-- An antimalarial DMARD used in SLE (first-line for skin/joint/mild disease, reduces flares and improves survival), rheumatoid arthritis, and antiphospholipid syndrome. Mechanisms: interference with antigen processing, Toll-like receptor inhibition, immunomodulation. Well tolerated; side effects: GI upset, rash, retinal toxicity (cumulative dose-related—require ophthalmologic screening), cardiomyopathy (rare), and myopathy. Monitor for retinopathy annually after 5 years or risk factors. Generally safe in pregnancy. Avoid with QT-prolonging drugs and in G6PD deficiency.

\medterm{IgG4-Related Disease}-- A fibro-inflammatory condition characterized by tumefactive lesions in multiple organs, dense lymphoplasmacytic infiltrate with >10 IgG4+ plasma cells/HPF and IgG4+/IgG+ ratio >40%, elevated serum IgG4 (>135 mg/dL), obliterative phlebitis, storiform fibrosis. Affects pancreas (type 1 autoimmune pancreatitis), salivary/lacrimal glands (Mikulicz disease), retroperitoneum, orbits, kidneys, lungs, aorta. Mimics malignancy, Sjogren, lymphoma. Treatment: glucocorticoids (first-line), rituximab, azathioprine, mycophenolate, methotrexate; disease-modifying anti-rheumatic drugs are often needed long-term.

\medterm{Iliotibial Band Syndrome}-- Lateral knee pain from IT band friction over lateral femoral epicondyle. Common in runners. Treatment: rest, stretching, strengthening, biomechanical correction.

\medterm{Immune-Mediated Necrotizing Myopathy}-- A subacute severe proximal myopathy with very high CK (often >10,000), minimal inflammation on biopsy, necrotic fibers. Two major autoantibodies: anti-HMGCR (3-hydroxy-3-methylglutaryl-CoA reductase, statin-associated or spontaneous) and anti-SRP (signal recognition particle). Anti-HMGCR: older, statin exposure, severe; anti-SRP: younger, more cardiac/respiratory involvement, worse prognosis. Treatment: high-dose glucocorticoids + rituximab or IVIG as first-line; additional immunosuppressants (methotrexate, azathioprine, mycophenolate) and physical therapy for rehabilitation.

\medterm{Inclusion Body Myositis}-- The most common acquired myopathy after age 50, characterized by asymmetric distal and proximal weakness (finger flexors, quadriceps), rimmed vacuoles, 15-18 nm tubulofilaments, TDP-43 pathology, endomysial inflammation with CD8+ T cells invading non-necrotic fibers. Refractory to immunotherapy. Clinical: dysphagia (40%), falls from quadriceps weakness. Diagnostic criteria (ENMC 2011): clinical + EMG + biopsy. Treatment: supportive (fall prevention, dysphagia management, speech therapy, and assistive devices for mobility. Clinical trials with arimoclomol, bimagrumab, and follistatin gene therapy are ongoing.

\medterm{Interferonopathies}-- Type I interferonopathies: Aicardi-Goutieres syndrome (TREX1, RNASEH2, SAMHD1 mutations, neonatal encephalopathy, calcifications, chilblains, SLE-like); STING-associated vasculopathy with onset in infancy (SAVI, TMEM173 mutations, severe interstitial lung disease, digital ischemia); Proteasome-associated autoinflammatory syndromes (PRAAS, PSMB8 mutations, CANDLE syndrome: lipodystrophy, fever, skin lesions, joint contractures). Treatment: JAK inhibitors (baricitinib, ruxolitinib) target IFN signalling. Baricitinib approved for SAVI and CANDLE. hematopoietic stem cell transplantation in severe Aicardi-Goutieres syndrome.

\medterm{Janus Kinase Inhibitors}-- JAK inhibitors (tofacitinib, baricitinib, upadacitinib, filgotinib, abrocitinib) block JAK-STAT signaling downstream of cytokine receptors (IL-6, IFN-gamma, IL-2, IL-7, IL-15, IL-21, GM-CSF, erythropoietin, thrombopoietin). Selectivity: JAK1/JAK3 (tofacitinib), JAK1/JAK2 (baricitinib), JAK1 (upadacitinib, filgotinib, abrocitinib). Indications: RA, PsA, axial SpA, UC, AD, alopecia areata. Safety: herpes zoster (vaccinate Shingrix before), VTE (risk factors: age, obesity, prior VTE), infections, lipid elevation, transaminitis, neutropenia, anemia, GI perforation. Drug interactions: CYP3A4 (tofacitinib, upadacitinib). Dose adjustment for renal impairment.

\medterm{Juvenile Arthritis}-- The most common chronic rheumatic disease of childhood (onset <16 years), encompassing 7 subtypes: systemic JIA (Still disease of childhood, quotidian fever, salmon rash, arthritis, hepatosplenomegaly, serositis, macrophage activation syndrome risk), oligoarticular JIA (4 or fewer joints in first 6 months, high risk of chronic anterior uveitis, ANA positive), polyarticular JIA (5 or more joints, RF-negative or RF-positive), psoriatic JIA (arthritis + psoriasis), enthesitis-related JIA (HLA-B27, male > female, enthesitis, acute uveitis), and undifferentiated JIA. See \hyperlink{Juvenile Idiopathic Arthritis}{juvenile idiopathic arthritis}.

\medterm{Juvenile Dermatomyositis}-- Juvenile dermatomyositis (JDM) presents with proximal weakness, heliotrope rash, Gottron papules, calcinosis (20-30%), lipodystrophy, vasculopathy (nailfold capillary changes). Anti-Mi-2 (classic DM, good prognosis), anti-NXP2 (calcinosis, malignancy risk low in children), anti-TIF1-gamma (malignancy risk in adults). Calcinosis is major morbidity. Treatment: glucocorticoids, methotrexate, hydroxychloroquine, IVIG for refractory; rituximab, TNF inhibitors, abatacept for refractory; bisphosphonates for calcinosis. JDM has better prognosis than adult DM; most children achieve remission but calcinosis can cause significant long-term disability.

\medterm{Juvenile Idiopathic Arthritis}-- The most common chronic rheumatic disease of childhood (onset <16 years), encompassing 7 subtypes: systemic JIA (Still disease of childhood, quotidian fever, salmon rash, arthritis, hepatosplenomegaly, serositis, macrophage activation syndrome risk), oligoarticular JIA (4 or fewer joints in first 6 months, high risk of chronic anterior uveitis, ANA positive), polyarticular JIA (5 or more joints, RF-negative or RF-positive), psoriatic JIA (arthritis + psoriasis), enthesitis-related JIA (HLA-B27, male > female, enthesitis, acute uveitis), and undifferentiated JIA. Treatment: NSAIDs, intra-articular corticosteroids, methotrexate, TNF inhibitors, IL-6 inhibitors (tocilizumab), IL-1 inhibitors (anakinra) for systemic JIA. Uveitis screening essential (ANA-positive, oligoarticular).


\medterm{Kawasaki Disease}-- Acute febrile vasculitis of childhood (peak 6 months-5 years), affecting medium-sized arteries, especially coronary arteries. Criteria: fever >=5 days plus 4 of 5: conjunctival injection, oral changes (strawberry tongue, cracked lips), cervical lymphadenopathy (>1.5 cm), polymorphous rash, extremity changes (erythema/edema of hands/feet, periungual desquamation). Treatment: IVIG (2 g/kg single infusion) + high-dose aspirin within 10 days of fever onset reduces coronary aneurysm risk from 25% to <5%. Recurrent Kawasaki disease possible.

\medterm{Lachman Test}-- ACL integrity test. Knee 20-30 degrees flexed, anterior drawer. Positive = ACL tear.

\medterm{Latarjet Procedure}-- Surgical procedure for shoulder instability using coracoid process transfer. Provides bony block to prevent recurrent dislocation.

\medterm{Lupus Erythematosus (Systemic Lupus Erythematosus - SLE)}-- A chronic multisystem autoimmune disease predominantly affecting women (9:1 female:male) of childbearing age, characterized by autoantibody production (ANA 98%, anti-dsDNA 70%, anti-Sm 30%) and immune complex deposition causing tissue injury. 2019 EULAR/ACR classification criteria: ANA positive at >=1:80 plus weighted clinical/immunologic domains (fever, hematologic, neuropsychiatric, mucocutaneous, serosal, musculoskeletal, renal). Clinical manifestations: mucocutaneous (malar rash, discoid rash, photosensitivity, oral ulcers, alopecia), serosal (pleurisy, pericarditis), musculoskeletal (arthritis, myositis), renal (lupus nephritis class I-VI on biopsy), neurologic (seizures, psychosis, mononeuritis multiplex, transverse myelitis), hematologic (hemolytic anemia, leukopenia, lymphopenia, thrombocytopenia, antiphospholipid syndrome), and pulmonary (interstitial lung disease, pulmonary hypertension). Lupus nephritis requires aggressive immunosuppression with mycophenolate or cyclophosphamide plus glucocorticoids. Hydroxychloroquine is recommended for all SLE patients. Belimumab (anti-BLyS), anifrolumab (anti-IFNAR), and calcineurin inhibitors (voclosporin) are newer targeted therapies. Pregnancy in SLE requires careful planning (avoid cyclophosphamide, mycophenolate), hydroxychloroquine continuation, and monitoring for flares and neonatal lupus.

\medterm{Lupus Nephritis}-- Renal involvement in SLE, classified into 6 classes by ISN/RPS: Class I (minimal mesangial), II (mesangial proliferative), III (focal proliferative, <50% glomeruli), IV (diffuse proliferative, >=50%—most severe), V (membranous), VI (advanced sclerosing). Treatment: Class III/IV: cyclophosphamide or mycophenolate + corticosteroids. Class V: mycophenolate or calcineurin inhibitor. Voclosporin approved for Class III/IV. Long-term monitoring essential.

\medterm{Macrophage Activation Syndrome}-- A life-threatening complication of rheumatic diseases (especially systemic JIA and SLE) characterized by uncontrolled cytokine release and hemophagocytosis. Criteria: fever, splenomegaly, cytopenias (>=2 lineages), hypertriglyceridemia or hypofibrinogenemia, hyperferritinemia (>684 mcg/L), elevated soluble CD25, low/absent NK cell activity. Treatment: high-dose corticosteroids, etoposide, anakinra, cyclosporine.

\medterm{Marfan Syndrome}-- Autosomal dominant connective tissue disorder (fibrillin-1 mutation). Features: tall stature, long limbs, arachnodactyly, pectus deformity, lens subluxation, aortic root dilation, mitral valve prolapse. Diagnosis: Ghent criteria. Management: beta-blockers, ARBs, elective aortic repair, regular cardiac monitoring.

\medterm{McMurray Test}-- Meniscal tear test. Knee flexed, rotated, extended. Click/pain indicates tear.

\medterm{Meniscal Tear}-- Tear of knee meniscus. Types: longitudinal, radial, flap, bucket handle, horizontal. Symptoms: knee pain, swelling, catching, locking. Diagnosis: clinical (McMurray, Thessaly), MRI. Treatment: conservative, arthroscopic partial meniscectomy, meniscal repair.

\medterm{Microscopic Polyangiitis}-- A necrotizing vasculitis affecting small blood vessels (capillaries, venules, arterioles) with few or no immune deposits (pauci-immune). Part of the ANCA-associated vasculitis group. Strongly associated with MPO-ANCA (p-ANCA, 80%). Presents with rapidly progressive glomerulonephritis (renal: hematuria, proteinuria, crescents), pulmonary capillaritis (hemoptysis, diffuse alveolar hemorrhage), palpable purpura, mononeuritis multiplex, and other organ involvement. Diagnosis: ANCA serology, renal or lung biopsy showing pauci-immune necrotizing inflammation. Treatment: high-dose corticosteroids + cyclophosphamide or rituximab for induction; azathioprine or rituximab for maintenance. Untreated, rapidly fatal.

\medterm{Mixed Connective Tissue Disease}-- An overlap syndrome with features of systemic lupus erythematosus, systemic sclerosis, and polymyositis/dermatomyositis. Highly associated with anti-U1 RNP antibodies. Presents with Raynaud phenomenon (nearly universal), puffy hands, arthritis, myositis, esophageal dysmotility, and serositis. Renal disease is less common than in SLE. Diagnosis: clinical + anti-U1 RNP positivity without sufficient criteria for other CTDs. Treatment: hydroxychloroquine, corticosteroids, and immunosuppressants (mycophenolate, methotrexate). Prognosis varies with predominant disease manifestation.


\medterm{Myositis}-- Inflammation of skeletal muscle. Includes polymyositis, dermatomyositis, inclusion body myositis, immune-mediated necrotizing myopathy. Diagnosis: CK elevation, EMG, muscle MRI, muscle biopsy. Treatment: immunosuppression with corticosteroids, methotrexate, azathioprine, mycophenolate, IVIG, rituximab.

\medterm{Nonsteroidal Anti-Inflammatory Drug}-- NSAID: drugs inhibiting cyclooxygenase (COX-1/COX-2), reducing prostaglandin synthesis—anti-inflammatory, analgesic, antipyretic. Types: nonselective (ibuprofen, naproxen, diclofenac, indomethacin, ketorolac) and COX-2 selective (celecoxib, etoricoxib). Uses: arthritis, musculoskeletal pain, gout, dysmenorrhea, postoperative pain. Adverse effects: GI ulceration/bleeding (reduce with PPI, COX-2), renal impairment (reduced GFR), cardiovascular risk (hypertension, fluid retention, increased MI risk—especially diclofenac), and asthma exacerbation. Avoid in renal disease, heart failure, GI bleeding, and late pregnancy; caution with anticoagulants. Use lowest effective dose/shortest duration.

\medterm{Osgood-Schlatter Disease}-- Traction apophysitis of tibial tubercle in adolescents. Pain/swelling at knee. Self-limiting, resolves with skeletal maturity. Treatment: rest, ice, NSAIDs.

\medterm{Osteoarthritis}-- The most common joint disease, characterized by progressive cartilage loss, subchondral bone sclerosis, osteophyte formation, and synovial inflammation, affecting knees, hips, hands (DIP/PIP Heberden/Bouchard nodes), spine. Risk factors: age, obesity, joint injury, malalignment, genetics, female sex. Clinical: activity-related pain, brief morning stiffness (<30 min), crepitus, bony enlargement, functional limitation. Diagnosis: clinical + radiographic (joint space narrowing, osteophytes, subchondral sclerosis, cysts). Management: non-pharmacological (weight loss, exercise, physiotherapy, joint protection), pharmacological (paracetamol, NSAIDs topical/oral, intra-articular corticosteroids, hyaluronic acid injections), and surgical (joint replacement) for end-stage disease.

\medterm{Osteogenesis Imperfecta}-- Hereditary disorder of collagen type I, causing brittle bones. Types I-VIII. Features: recurrent fractures, blue sclerae, hearing loss, dentinogenesis imperfecta, short stature. Diagnosis: clinical, genetic testing. Management: bisphosphonates, orthopedic surgery, physical therapy.

\medterm{Osteomalacia}-- Softening of bones due to defective mineralization, usually from vitamin D deficiency. Adults; rickets in children. Features: bone pain, muscle weakness, fractures. Diagnosis: low calcium/phosphate, high alkaline phosphatase, low vitamin D, radiographic changes (Looser zones). Treatment: vitamin D supplementation, calcium.

\medterm{Osteomyelitis}-- Bone infection, usually bacterial (Staphylococcus aureus). Acute: fever, pain, swelling. Chronic: sinus tracts, sequestrum, involucrum. Diagnosis: MRI, bone scan, biopsy. Treatment: antibiotics (prolonged), surgical debridement.

\medterm{Osteonecrosis (Avascular Necrosis)}-- Osteonecrosis (AVN) is bone death from vascular compromise, affecting femoral head (most common), humeral head, femoral condyles, talus. Risk factors: glucocorticoids (high dose, cumulative), alcohol, trauma, sickle cell disease, SLE, antiphospholipid syndrome, pancreatitis, decompression sickness, radiation, organ transplantation. ARCO classification: stage 0 (normal X-ray, MRI+), I (normal X-ray, MRI+), II (crescent sign), III (subchondral collapse), IV (joint space narrowing, OA). Treatment: core decompression for early stages, osteotomy, and joint replacement for advanced collapse. Bisphosphonates may reduce risk in high-dose glucocorticoid therapy.

\medterm{Overlap Myositis}-- Patients with features of inflammatory myopathy (DM/PM/IMNM) plus another systemic autoimmune disease (SLE, SSc, MCTD, Sjogren). Anti-PM-Scl (Scl/PM overlap, limited SSc + myositis, good prognosis); anti-Ku (SLE/PM overlap); anti-U1-RNP (MCTD with myositis). Treatment: tailored to predominant organ involvement; often requires combination immunosuppression.

\medterm{Paget Disease of Bone}-- Chronic bone disorder with excessive bone breakdown and formation, leading to enlarged/deformed bones. Features: bone pain, deformity, fractures, hearing loss (skull involvement), high alkaline phosphatase. Diagnosis: radiographic (cotton wool skull, flame-shaped lytic lesions), bone scan. Treatment: bisphosphonates (alendronate, zoledronic acid), calcitonin.

\medterm{Palindromic Rheumatism}-- Characterized by recurrent, self-limited attacks of mono- or oligoarthritis (hours to days), periarticular soft tissue swelling, no radiographic erosions between attacks. 30-50% evolve to RA (anti-CCP positive predicts progression). Treatment: NSAIDs for attacks; hydroxychloroquine reduces frequency; low-dose colchicine; DMARDs if progressing to RA.

\medterm{Patellofemoral Pain Syndrome}-- Anterior knee pain from patellar maltracking. Common in runners, cyclists. Symptoms: pain with stairs, squatting, prolonged sitting. Treatment: quadriceps strengthening, patellar taping, orthotics.

\medterm{Phalen Test}-- Wrist flexion test for carpal tunnel syndrome. Symptoms reproduced within 60 seconds.

\medterm{Piriformis Syndrome}-- Sciatic nerve compression by piriformis muscle. Buttock pain, radiating down leg. Diagnosis: clinical, exclude disc herniation. Treatment: stretching, injection, surgical release.

\medterm{Plantar Fasciitis}-- Inflammation of plantar fascia causing heel pain, worse with first steps in morning. Risk factors: obesity, tight calf muscles, prolonged standing. Treatment: stretching, orthotics, night splints, corticosteroid injection, shockwave therapy, surgery if refractory.

\medterm{Polyarteritis Nodosa}-- Necrotizing inflammation of medium-sized muscular arteries, spares lungs. Associated with hepatitis B (10-30%). Clinical: systemic symptoms, mononeuritis multiplex, livedo reticularis, renal involvement (hypertension, renal insufficiency—no glomerulonephritis), GI ischemia. Diagnosis: angiography (microaneurysms, "string of beads"), biopsy. Treatment: corticosteroids + cyclophosphamide for severe disease; antiviral for HBV-associated.

\medterm{Polymyalgia Rheumatica}-- An inflammatory syndrome of older adults (age >50, peak 70-80) causing bilateral shoulder and pelvic girdle pain and stiffness lasting >45 minutes, with elevated inflammatory markers (ESR >40, CRP elevated). Strongly associated with giant cell arteritis (15-20% develop GCA; 40-50% of GCA patients have PMR). 2012 ACR/EULAR provisional classification: age >=50, bilateral shoulder pain, abnormal ESR/CRP, morning stiffness >45 min, new hip pain, absence of RF/anti-CCP, absence of other inflammatory arthritis. Treatment: low-dose prednisolone 15-25 mg/day (not exceeding 20 mg/day in most guidelines) with rapid taper; methotrexate as steroid-sparing; tocilizumab for refractory cases. Relapse common during steroid taper.


\medterm{Polymyositis}-- An idiopathic inflammatory myopathy characterized by symmetric, proximal muscle weakness (shoulders, hips) with elevated muscle enzymes (CK, aldolase, LDH). More common in women, peaks at 30-50 years. Diagnosis: EMG (myopathic units with fibrillations), muscle MRI (edema), and muscle biopsy (CD8+ T cell infiltration of MHC-I expressing muscle fibers). Autoantibodies: anti-Jo-1 (antisynthetase syndrome with interstitial lung disease, arthritis, mechanic's hands). Treatment: high-dose corticosteroids, methotrexate, azathioprine, mycophenolate, and IVIG for refractory disease.

\medterm{Posterior Cruciate Ligament Tear}-- PCL tear: less common, usually from direct blow to tibia (dashboard injury). Symptoms: knee pain, swelling, instability. Diagnosis: posterior drawer test, MRI. Treatment: usually conservative, surgical if combined injuries.

\medterm{Posterior Drawer Test}-- PCL integrity test. Knee flexed 90 degrees, posterior tibial translation.

\medterm{Pseudogout}-- Calcium pyrophosphate deposition disease (CPPD): deposition of calcium pyrophosphate dihydrate crystals in joints, causing acute arthritis resembling gout (pseudogout) or chronic osteoarthritis-like disease. Affects knees, wrists, and metacarpophalangeal joints; associated with aging, hemochromatosis, hyperparathyroidism, hypomagnesemia, hypothyroidism. Acute attacks: swelling, pain, effusion with weakly positively birefringent (rhomboid) crystals on polarized microscopy, chondrocalcinosis on X-ray. Management: NSAIDs, colchicine, intra-articular corticosteroids for acute; treat associated metabolic disease. Distinguish from gout (monosodium urate, negatively birefringent).

\medterm{Psoriatic Arthritis}-- An inflammatory arthritis associated with psoriasis (precedes arthritis in 70%, simultaneous 15%, follows 15%), affecting 20-30% of psoriasis patients. CASPAR criteria: inflammatory articular disease plus psoriasis (2 points), nail dystrophy (1), negative RF (1), dactylitis (1), juxta-articular new bone formation (1); score >=3. Patterns: distal interphalangeal predominant, asymmetric oligoarthritis, symmetric polyarthritis (RA-like), spondylitis, arthritis mutilans. Treatment: NSAIDs for mild disease; methotrexate, leflunomide, sulfasalazine as first-line DMARDs; TNF inhibitors, IL-17 inhibitors (secukinumab, ixekizumab), IL-23 inhibitors (guselkumab), and JAK inhibitors (tofacitinib, upadacitinib) for moderate-to-severe disease; physiotherapy and occupational therapy.


\medterm{Pyoderma Gangrenosum}-- A neutrophilic dermatosis causing rapidly expanding, painful ulcers with violaceous undermined borders, pathergy phenomenon. Associated with IBD (50%), hematologic malignancy, rheumatoid arthritis, PAPA syndrome. Peristomal PG in ostomy patients. Treatment: high-potency topical steroids, intralesional steroids, systemic corticosteroids, cyclosporine, TNF inhibitors (infliximab, adalimumab); avoid surgical debridement (worsens via pathergy).

\medterm{Raynaud Syndrome}-- A vasospastic disorder causing episodic digital color changes (white -> blue -> red) triggered by cold or stress, due to exaggerated vasoconstriction of digital arteries. Primary Raynaud: benign, symmetric, in young women, no underlying disease. Secondary Raynaud: associated with scleroderma, SLE, mixed connective tissue disease, vasculitis, and cryoglobulins—often asymmetric, severe, with digital ulcers/gangrene. Diagnosis: clinical; nailfold capillaroscopy and autoantibody testing to detect secondary causes. Management: avoid cold, smoking cessation, calcium channel blockers (nifedipine), prostacyclin analogs (iloprost) for severe; treat underlying disease; avoid sympathomimetics.

\medterm{Reactive Arthritis}-- A post-infectious spondyloarthritis occurring 1-4 weeks after genitourinary (Chlamydia trachomatis) or gastrointestinal (Salmonella, Shigella, Campylobacter, Yersinia) infection. Classic triad (Reiter syndrome): arthritis, conjunctivitis/uveitis, urethritis/cervicitis (only 1/3 have full triad). HLA-B27 positive in 60-80%. Clinical: asymmetric oligoarthritis (lower extremities), enthesitis (Achilles, plantar fascia), dactylitis, circinate balanitis, keratoderma blennorrhagicum (palmar/plantar pustulosis). Treatment: NSAIDs, corticosteroids (oral or intra-articular), DMARDs (sulfasalazine, methotrexate) for persistent arthritis, TNF inhibitors for refractory cases. Most cases resolve over months but may recur with subsequent infections.

\medterm{Remitting Seronegative Symmetrical Synovitis with Pitting Edema}-- RS3PE: an acute-onset, symmetric polyarthritis (usually in elderly) with pitting edema of the hands/feet, negative rheumatoid factor and anti-CCP, and marked response to low-dose corticosteroids. Distinct from RA—tends to be severe but responds dramatically to steroids and often remits; some cases progress to RA or other rheumatic disease (polymyalgia rheumatica, spondyloarthritis, or paraneoplastic). Features: symmetric tenosynovitis of wrist/hands with edema. Management: corticosteroids (rapid response), NSAIDs; assess for underlying malignancy or connective tissue disease.

\medterm{Rheumatoid Arthritis}-- A chronic, symmetric, erosive inflammatory polyarthritis primarily affecting small joints of the hands and feet, with a female-to-male ratio of 3:1 and peak onset 40-60 years. Pathogenesis involves genetic susceptibility (HLA-DRB1 shared epitope), environmental triggers (smoking, periodontal disease), and immune dysregulation (citrullination of proteins, anti-CCP antibody production, rheumatoid factor, cytokine-driven synovial inflammation with TNF, IL-6, IL-1). Clinical: symmetric small joint polyarthritis (MCP, PIP, MTP), morning stiffness >1 hour, rheumatoid nodules, extra-articular manifestations (Sjogren, interstitial lung disease, pericarditis, vasculitis, Felty syndrome). Diagnosis: 2010 ACR/EULAR criteria: joint involvement (0-5), serology (RF and anti-CCP, 0-3), acute phase reactants (0-1), symptom duration (0-1). Treatment: treat-to-target strategy with methotrexate as anchor DMARD, plus short-term glucocorticoids; early use of TNF inhibitors, IL-6 inhibitors (tocilizumab), abatacept, JAK inhibitors (tofacitinib, baricitinib, upadacitinib) for inadequate response; rituximab for seropositive disease. Intensive disease management reduces joint damage and disability.

\medterm{Rheumatoid Factor}-- RF: an autoantibody (usually IgM) directed against the Fc portion of IgG, positive in ~70-80 percent of rheumatoid arthritis (seropositive RA) but also in other conditions: Sjogren, SLE, chronic infection (TB, endocarditis, hepatitis C), cryoglobulinemia, sarcoidosis, and healthy elderly. High titers correlate with more severe, erosive RA. Used with anti-CCP (more specific for RA) for diagnosis and classification (ACR/EULAR criteria). Serial monitoring not generally recommended. Negative RF does not exclude RA (seronegative RA).

\medterm{Rickets}-- See osteomalacia. Pediatric form with growth plate changes, bowing of legs, craniotabes, rachitic rosary.

\medterm{Sarcoidosis}-- A multisystem granulomatous disorder of unknown etiology characterized by non-caseating granulomas, predominantly affecting lungs (90%), lymph nodes, skin, eyes, liver, heart, nervous system. Pulmonary stages (Scadding): 0 (normal), I (bilateral hilar lymphadenopathy), II (BHL + parenchymal), III (parenchymal only), IV (fibrosis). Extrapulmonary: Lofgren syndrome (acute, erythema nodosum, bilateral hilar adenopathy, arthritis, fever - good prognosis), Heerfordt syndrome (parotid enlargement, uveitis, fever, facial nerve palsy), cardiac sarcoid (conduction abnormalities, cardiomyopathy, sudden death), neurosarcoid (cranial nerve palsies, meningitis, brain parenchymal lesions), cutaneous (lupus pernio, erythema nodosum). Diagnosis: biopsy of affected tissue showing non-caseating granulomas, exclusion of other granulomatous diseases (TB, fungal). Treatment: corticosteroids first-line; steroid-sparing (methotrexate, azathioprine, mycophenolate, leflunomide); TNF inhibitors (infliximab, adalimumab) for refractory disease.

\medterm{Scleroderma}-- Systemic sclerosis: an autoimmune connective tissue disease of skin and internal organs due to fibrosis, vasculopathy, and autoantibodies. Limited cutaneous (CREST—calcinosis, Raynaud, esophageal dysmotility, sclerodactyly, telangiectasia) and diffuse cutaneous forms. Features: skin thickening (sclerodactyly, skin tightness), Raynaud phenomenon, digital ulcers, pulmonary fibrosis, pulmonary hypertension, renal crisis (hypertension, renal failure—emergency), GI dysmotility. Diagnosis: clinical, ANA, anti-Scl-70 (diffuse), anti-centromere (limited), anti-RNA polymerase III. Management: immunosuppression (mycophenolate, methotrexate, cyclophosphamide), vasodilators (pulmonary hypertension—endothelin antagonists, PDE5 inhibitors), ACE inhibitors for renal crisis, symptom control.

\medterm{Septic Arthritis}-- Joint infection, usually bacterial (Staph, Strept, Gonorrhea). Emergency: fever, severe joint pain, swelling, inability to bear weight. Diagnosis: joint aspiration (WBC >50,000, positive culture). Treatment: antibiotics, joint drainage (arthroscopy/open).

\medterm{Seronegative Arthritis}-- Arthritis with negative rheumatoid factor and anti-CCP. Includes spondyloarthropathies, psoriatic arthritis, reactive arthritis.

\medterm{Sever Disease}-- Calcaneal apophysitis in children/adolescents. Heel pain, worse with activity. Treatment: rest, heel cups, stretching.

\medterm{Shoulder Impingement}-- Compression of rotator cuff tendons beneath acromion. Symptoms: shoulder pain with overhead activity, weakness. Diagnosis: clinical (Neer, Hawkins tests), ultrasound/MRI. Treatment: physical therapy, corticosteroid injection, surgical decompression.

\medterm{Sicca Syndrome}-- Dry eyes and dry mouth due to decreased tear and saliva production. Primary: Sjogren syndrome. Secondary: associated with RA, SLE, scleroderma, graft-versus-host disease, hepatitis C, sarcoidosis, HIV, medications (anticholinergics, antidepressants). Diagnosis: Schirmer test (<5 mm/5 min), rose bengal staining, salivary gland biopsy. Treatment: artificial tears, pilocarpine, cevimeline, dental care.

\medterm{Sinding-Larsen-Johansson Syndrome}-- Traction apophysitis of inferior patellar pole in adolescents. Similar to Osgood-Schlatter but at knee. Treatment: rest, activity modification.

\medterm{Sjogren Syndrome}-- Primary Sjogren syndrome (pSS) is a chronic autoimmune exocrinopathy characterized by lymphocytic infiltration of lacrimal and salivary glands causing dry eyes (keratoconjunctivitis sicca) and dry mouth (xerostomia), with female predominance (9:1) and peak onset 40-60. Secondary SS occurs with RA, SLE, SSc. 2016 ACR/EULAR criteria: anti-Ro/SSA positive (3 points) OR focal lymphocytic sialadenitis with focus score >=1 (3 points) OR ocular staining score >=5 (1 point) OR Schirmer test <=5 mm/5 min (1 point). Treatment: symptomatic (artificial tears, saliva substitutes, pilocarpine, cevimeline); systemic immunosuppression (hydroxychloroquine, methotrexate, azathioprine, mycophenolate, rituximab) for extraglandular manifestations (arthritis, ILD, vasculitis, nephritis, neuropathy). Lymphoma risk (3-5 percent, MALT in parotid) requires long-term surveillance.

\medterm{SLAP Lesion}-- Superior labrum anterior to posterior tear. Associated with overhead activities, trauma. Treatment: rest, physical therapy, surgical repair if symptomatic.

\medterm{Snapping Hip Syndrome}-- Hip pain/clicking from tendon snapping over bony prominence. Internal (iliopsoas), external (IT band), intra-articular (labral tear). Treatment: physical therapy, injection, surgery if refractory.

\medterm{Soft Tissue Rheumatism}-- Pain conditions affecting muscles, tendons, ligaments, bursae. Examples: frozen shoulder, lateral epicondylitis, rotator cuff tendinopathy. Treatment: rest, physiotherapy, injections, surgery if refractory.

\medterm{Spondyloarthritis}-- Group of inflammatory arthritis associated with HLA-B27. Includes ankylosing spondylitis, psoriatic arthritis, reactive arthritis, enteropathic arthritis. Features: axial involvement, enthesitis, dactylitis.

\medterm{Sulfasalazine}-- A DMARD combining sulfapyridine and 5-aminosalicylic acid, used in rheumatoid arthritis, psoriatic arthritis, and inflammatory bowel disease (ulcerative colitis). Slows disease progression in RA; onset over 1-3 months. Adverse effects: GI upset, rash, headache, myelosuppression (monitor CBC), hepatotoxicity, male infertility (reversible), and sulfa allergy. Contraindicated in sulfonamide allergy and G6PD deficiency; use folate supplementation. Useful in combination DMARD therapy and for patients intolerant to methotrexate.

\medterm{Sweet Syndrome}-- Sweet syndrome (acute febrile neutrophilic dermatosis) presents with fever, neutrophilia, and tender erythematous plaques/nodules with dense dermal neutrophilic infiltrate. Classical (idiopathic), malignancy-associated (AML, solid tumors), drug-induced (G-CSF, granulocyte colony-stimulating factor). Associated with VEXAS syndrome. Treatment: systemic corticosteroids (rapid response); colchicine, dapsone, indomethacin as steroid-sparing; treat underlying malignancy.

\medterm{Synovitis}-- Inflammation of the synovial membrane lining a joint, causing swelling, pain, warmth, and joint effusion. Causes: inflammatory arthritis (RA, psoriatic, gout, reactive, septic), autoimmune (SLE), and trauma/overuse. Diagnosis: clinical (joint swelling, effusion), ultrasound (synovial hypertrophy, Doppler signal), MRI, and joint aspiration for septic vs. crystal arthritis (gout—urate crystals, pseudogout—CPPD). Chronic synovitis causes cartilage and bone destruction. Management: treat underlying cause—NSAIDs, corticosteroids, DMARDs/biologics for inflammatory arthritis, antibiotics for septic arthritis (emergency—urgent aspiration and antibiotics), and intra-articular steroid injection.

\medterm{Systemic Juvenile Idiopathic Arthritis}-- sJIA: a subtype of juvenile idiopathic arthritis with prominent systemic inflammation—daily spiking fevers, evanescent salmon-pink rash, arthritis, lymphadenopathy, hepatosplenomegaly, and elevated inflammatory markers (ferritin). Complications: macrophage activation syndrome (MAS—emergency with cytopenias, coagulopathy, very high ferritin), growth impairment, uveitis. Diagnosis: clinical; exclude infection/malignancy. Management: NSAIDs, corticosteroids, methotrexate, IL-1 inhibitors (anakinra, canakinumab), IL-6 inhibitor (tocilizumab); early aggressive treatment improves outcomes.

\medterm{Systemic Lupus Erythematosus}-- Multisystem autoimmune disease, more common in women. Features: malar rash, photosensitivity, oral ulcers, arthritis, serositis, renal disorder, neurologic disorder, hematologic disorder, immunologic disorder. Diagnosis: ANA positive plus clinical/immunologic criteria. Treatment: hydroxychloroquine, corticosteroids, immunosuppressants (mycophenolate, cyclophosphamide, azathioprine), biologics (belimumab, anifrolumab).

\medterm{Systemic Sclerosis}-- Systemic sclerosis: an autoimmune connective tissue disease of skin and internal organs due to fibrosis, vasculopathy, and autoantibodies. See \hyperlink{Scleroderma}{scleroderma}.

\medterm{Takayasu Arteritis}-- Large-vessel vasculitis affecting the aorta and its major branches, predominantly in young Asian women (<40 years). Early stage: constitutional symptoms, arthralgia. Late stage: pulselessness, blood pressure discrepancy (>10 mmHg between arms), claudication, bruits, hypertension (renal artery stenosis), visual symptoms. Angiography shows wall thickening, stenosis, aneurysms. Treatment: corticosteroids, steroid-sparing agents (methotrexate, azathioprine, mycophenolate), TNF inhibitors; angioplasty/stenting for severe stenosis.

\medterm{Tarsal Tunnel Syndrome}-- Tibial nerve compression at ankle. Symptoms: burning pain, tingling in foot sole. Tinel sign at ankle. Treatment: orthotics, corticosteroid injection, surgical decompression.

\medterm{Tendinitis}-- Inflammation of tendon. Common sites: rotator cuff, patellar, Achilles, lateral epicondyle (tennis elbow), medial epicondyle (golfer's elbow). Treatment: rest, ice, NSAIDs, physical therapy, corticosteroid injection, surgery if refractory.

\medterm{Tendonitis}-- Tendinopathy: inflammation or degeneration of a tendon from overuse, repetitive strain, or trauma, causing localized pain, swelling, and tenderness with activity. Common sites: rotator cuff (supraspinatus), lateral elbow (tennis elbow), medial elbow (golfer's elbow), Achilles, patellar (jumper's knee), biceps, and De Quervain. Diagnosis: clinical; ultrasound/MRI for severity or rupture. Management: relative rest, ice, NSAIDs, physiotherapy (eccentric exercises), corticosteroid injection (caution—may weaken tendon), platelet-rich plasma; surgical repair for rupture (e.g., Achilles). Avoid prolonged immobilization; gradual return to activity.

\medterm{Tenosynovitis}-- Inflammation of a tendon and its synovial sheath, causing pain, swelling, and crepitus along the tendon. Causes: overuse/repetitive activity (De Quervain—first dorsal compartment, trigger finger), infection (septic tenosynovitis—emergency, especially flexor sheaths of hand, gonococcal), inflammatory arthritis (RA, psoriatic, gout, reactive arthritis), and trauma. Features: tenderness along tendon, swelling, pain with movement, and in infection: fusiform swelling, pain on passive extension, tenderness along sheath. Diagnosis: clinical; ultrasound/MRI; culture for infection. Management: rest, NSAIDs, splinting, corticosteroid injection, treat underlying disease; surgical drainage for septic tenosynovitis.

\medterm{Thoracic Outlet Syndrome}-- Compression of neurovascular bundle at thoracic outlet. Types: neurogenic (most common), venous, arterial. Symptoms: neck/shoulder/arm pain, paresthesia, weakness, vascular symptoms. Diagnosis: clinical, imaging, nerve conduction. Treatment: physical therapy, surgical decompression.

\medterm{Tinel Sign}-- Tapping over nerve produces tingling in distribution. Indicates nerve irritation/regeneration. Used for carpal tunnel, cubital tunnel, tarsal tunnel.

\medterm{Tophus}-- Deposits of monosodium urate crystals in soft tissues, characteristic of chronic gout. Appear as subcutaneous nodules over joints, ears, olecranon, Achilles tendon. Can erode through skin, forming draining sinus tracts. Radiographically: punched-out erosions with overhanging edges. Treatment with urate-lowering therapy causes tophi to dissolve over months to years.

\medterm{Trigger Finger}-- Stenosing tenosynovitis of flexor tendon, causing locking/catching. Treatment: splinting, corticosteroid injection, surgical release.

\medterm{Tumor Necrosis Factor Inhibitor}-- TNF inhibitor (biologic): monoclonal antibodies/fusion proteins blocking TNF-alpha—infliximab, adalimumab (monoclonals), etanercept (soluble receptor), certolizumab, golimumab. Used in rheumatoid arthritis (after methotrexate failure), ankylosing spondylitis, psoriatic arthritis, psoriasis, IBD (Crohn, UC). Advantages: rapid, effective control of inflammation and joint damage. Risks: infection (tuberculosis—screen before starting, hepatitis B reactivation, fungal), lymphoma (rare), demyelination, heart failure worsening, and injection/infusion reactions. Screening (TB, HBV) and live vaccine avoidance required; contraindicated in active infection.

\medterm{Valgus/Varus Stress Test}-- Collateral ligament integrity. Valgus (MCL), varus (LCL).

\medterm{Vasculitis}-- Inflammation of blood vessels causing tissue ischemia and necrosis. Classification by vessel size. Large vessel: giant cell arteritis, Takayasu arteritis. Medium vessel: polyarteritis nodosa (PAN), Kawasaki disease. Small vessel: ANCA-associated vasculitis (microscopic polyangiitis, granulomatosis with polyangiitis, eosinophilic granulomatosis with polyangiitis), IgA vasculitis (Henoch-Schonlein purpura), and anti-GBM disease. Diagnosis: clinical presentation, ANCA, ANA, complement, and tissue biopsy. Treatment: corticosteroids and immunosuppressants (cyclophosphamide, rituximab for ANCA-associated).

\medterm{Vasculitis (Systemic)}-- See vasculitis. Large vessel (GCA, Takayasu), medium vessel (PAN, Kawasaki), small vessel (ANCA-associated, IgA vasculitis).

\medterm{VEXAS Syndrome}-- VEXAS (Vacuoles, E1 ubiquitin-activating enzyme, X-linked, Autoinflammatory, Somatic) syndrome is an adult-onset autoinflammatory disorder caused by somatic UBA1 mutations in hematopoietic stem cells, affecting men >50. Clinical: recurrent fever, chondritis (relapsing polychondritis-like), vasculitis (Sweet-like, polyarteritis-like), macrocytic anemia, thrombocytopenia, vacuoles in myeloid/erythroid precursors. High mortality from infections, thrombosis, hematologic progression (MDS, AML). Treatment: high-dose corticosteroids for acute inflammation; JAK inhibitors (ruxolitinib, tofacitinib); azacitidine for MDS; supportive care for cytopenias. hematopoietic stem cell transplantation is the only potential cure.

\medterm{Synovial Fluid Analysis}-- Laboratory examination of aspirated joint fluid for appearance, viscosity, leukocyte count and differential, crystals under polarized light, Gram stain, and culture. It helps distinguish noninflammatory, inflammatory, crystal-associated, and septic arthritis.

\medterm{Arthrocentesis}-- Aspiration of synovial fluid from a joint using a needle for diagnosis or symptom relief. Sterile technique is essential, and suspected septic arthritis requires prompt fluid analysis and culture before or alongside treatment.

\medterm{C-Reactive Protein (CRP)}-- An acute-phase protein synthesized by the liver in response primarily to interleukin-6 and other inflammatory signals. It is useful for assessing inflammatory activity and response to treatment but is nonspecific and should not be interpreted in isolation.


\medterm{Treat-to-Target Strategy}-- A management approach in which treatment is adjusted at regular intervals toward a predefined goal such as remission or low disease activity. It requires validated disease-activity measures, shared decisions, monitoring, and timely escalation or modification.

\medterm{Disease-Modifying Antirheumatic Drug (DMARD)}-- A medication that reduces inflammatory disease activity and helps prevent structural damage or disability rather than providing only symptomatic relief. Conventional synthetic, biologic, and targeted synthetic DMARDs are used according to disease and patient factors.

\medterm{Glucocorticoid-Sparing Therapy}-- Use of a non-glucocorticoid treatment strategy to control inflammation while reducing cumulative corticosteroid exposure. It aims to limit osteoporosis, infection, diabetes, hypertension, adrenal suppression, and other glucocorticoid complications.

\medterm{Gout Flare}-- An acute inflammatory response to monosodium urate crystal deposition in or around a joint. It commonly causes abrupt severe pain, swelling, warmth, and erythema, often involving the first metatarsophalangeal joint; aspiration can confirm crystals and exclude infection.

\medterm{Urate-Lowering Therapy}-- Long-term treatment that reduces serum urate to dissolve monosodium urate deposits and prevent gout flares and structural damage. Xanthine oxidase inhibitors and uricosuric medicines are selected according to renal function, comorbidity, and treatment goals.

\medterm{Pseudogout Flare}-- Acute crystal-induced arthritis caused by calcium pyrophosphate dihydrate deposition, often affecting the knee or wrist. Identification of positively birefringent rhomboid crystals in synovial fluid supports the diagnosis.


\medterm{Giant Cell Arteritis Emergency}-- Suspected giant cell arteritis with visual symptoms, jaw or tongue claudication, or neurologic ischemia requires immediate glucocorticoid treatment to reduce risk of irreversible visual loss while diagnostic evaluation proceeds.

\medterm{Immunosuppression Screening}-- Pre-treatment evaluation for latent or active infection, vaccination status, malignancy risk, and organ dysfunction before immunosuppressive therapy. Testing commonly includes tuberculosis, hepatitis B and C, HIV, blood counts, liver and kidney function, and pregnancy when relevant.

\medterm{Raynaud Phenomenon}-- Episodic vasospasm of digital arteries triggered by cold or emotional stress, producing sequential color changes and altered sensation. Primary disease is usually benign; secondary Raynaud phenomenon may signal systemic sclerosis, lupus, vasculitis, or another connective-tissue disease.

\medterm{Inflammatory Back Pain}-- A pattern of chronic back pain characterized by insidious onset, improvement with exercise, lack of improvement with rest, nocturnal pain, and prolonged morning stiffness. It suggests axial spondyloarthritis but is not diagnostic without supportive clinical, laboratory, and imaging findings.
\medterm{Cryoglobulinemic vasculitis}-- Systemic vasculitis caused by cryoglobulin-containing immune complexes that deposit in small and medium vessels, classically presenting with purpura, arthralgia, and glomerulonephritis. Most cases are associated with hepatitis C virus infection.
\medterm{Cutaneous lupus}-- A form of lupus erythematosus limited to the skin, manifesting as discoid lesions, subacute cutaneous plaques, or acute malar erythema. It may occur in isolation or as a feature of systemic lupus erythematosus.
\medterm{Jaccoud's arthropathy}-- A non-erosive, reducible deformity of the metacarpophalangeal joints caused by chronic periarticular fibrosis and tendon laxity, most commonly associated with systemic lupus erythematosus.
\medterm{Lupus pernio}-- A chronic, violaceous skin lesion of the face and extremities associated with sarcoidosis, indicating a higher risk of systemic involvement and chronic disease.
\medterm{Polychondritis}-- Relapsing inflammatory disorder targeting cartilage throughout the body, causing pain, swelling, and destruction of the ears, nose, larynx, trachea, and joints. May lead to airway collapse and hearing loss.
\medterm{Primary hypertrophic osteoarthropathy}-- A rare inherited syndrome characterized by digital clubbing, periostosis of long bones, and pachydermia, presenting with joint pain, skin thickening, and visible nail-bed enlargement.
\medterm{Anti-citrullinated peptide antibody}-- An autoantibody highly specific for rheumatoid arthritis that may precede clinical disease and is associated with erosive joint damage.

\medterm{Sacroiliitis}-- Inflammation of one or both sacroiliac joints, causing pain in the lower back, buttock, or posterior thigh that may worsen with prolonged rest. Causes include axial spondyloarthritis, psoriatic or enteropathic arthritis, infection, trauma, and degenerative disease. Diagnosis combines clinical assessment with inflammatory markers and sacroiliac imaging, particularly magnetic resonance imaging when early inflammation is suspected.
